\beforesec
\section{Weight-space ensembles for fine-tuning} 
\label{sec:main}
\postsec
This section describes and motivates our proposed method, WiSE-FT, which consists of two simple steps.
First, we fine-tune the zero-shot model on application-specific data. 
Second, we combine the original zero-shot and fine-tuned models by linearly interpolating between their weights, also referred to as weight-space ensembling.
WiSE-FT can be implemented in a few lines of PyTorch, and we provide example code in Appendix \ref{sec:pseudo-code}.

The zero-shot model excels under distribution shift while standard fine-tuning achieves high accuracy on the reference distribution.
Our motivation is to combine these two models into one that achieves the best of both worlds.
Weight-space ensembles are a natural choice as they ensemble without extra computational cost. 
Moreover,
previous work has suggested that interpolation in weight space
may improve performance when models share part of their optimization trajectory \cite{izmailov2018averaging, neyshabur2020being}.

\smallpara{Step 1: Standard fine-tuning.} 
As in Section~\ref{sec:expsetup}, we let $\mathcal{S}_{\mathrm{ref}}^\text{tr}$ denote the dataset used for fine-tuning and $g$ denote the image encoder used by CLIP.
We are now explicit in writing $g\mleft(x, \mathbf{V}_{\text{enc}} \mright)$ where $x$ is an input image and $\mathbf{V}_{\text{enc}}$ are the parameters of the encoder $g$.
Standard fine-tuning considers the model $f\mleft(x, \theta\mright) = g\left(x, \mathbf{V}_{\text{enc}}\right)^\top \textbf{W}_{\text{classifier}}$ where $\textbf{W}_{\text{classifier}} \in \mathbb{R}^{d \times k}$ is the classification head and $\theta = \left[\mathbf{V}_{\text{enc}}, \textbf{W}_{\text{classifier}} \right]$ are the parameters of $f$. 
We then solve $\argmin_{\theta} \left\{ \sum_{(x_i, y_i) \in \mathcal{S}^\text{tr}_\mathrm{ref}} \ell\mleft(f(x_i,\theta),\ y_i \mright) + \lambda R(\theta) \right\}$ where $\ell$ is the cross-entropy loss and $R$ is a regularization term (e.g., weight decay).
We consider the two most common variants of fine-tuning: end-to-end, where all values of $\theta$ are modified, and fine-tuning only a linear classifier, where $\mathbf{V}_{\text{enc}}$ is fixed at the value learned during pre-training. 
Appendices~\ref{sec:e2e-ft} and \ref{sec:moreclf} provide additional details. %

\smallpara{Step 2: Weight-space ensembling.}
For a \textit{\ALPHA{}}  $\alpha \in [0,1]$, we consider the \textit{weight-space ensemble} between the zero-shot model with parameters $\theta_0$ and the model obtained via standard fine-tuning with parameters $\theta_1$. 
The predictions of the weight-space ensemble $\mathsf{wse}$ are given by
\begin{align}\label{eqn:wse}
    \mathsf{wse}(x, \alpha) = f\mleft(x, (1-\alpha)\cdot\theta_0 + \alpha\cdot\theta_1 \mright)\ ,
\end{align}
i.e., we use the element-wise weighted average of the zero-shot and fined-tuned parameters. 
When fine-tuning only the linear classifier, weight-space ensembling is equivalent to the traditional output-space ensemble \cite{dietterich2000ensemble, breiman1996bagging, FREUND1997119} $(1-\alpha)\cdot f\mleft(x, \theta_0\mright)  + \alpha\cdot f\mleft(x, \theta_1 \mright)$ since Equation~\ref{eqn:wse} decomposes as $(1-\alpha)\cdot g\mleft(x, \mathbf{V}_{\text{enc}}\mright)^\top \textbf{W}_{\text{zero-shot}} + \alpha \cdot g\mleft(x, \mathbf{V}_{\text{enc}}\mright)^\top \textbf{W}_{\text{classifier}}$.
 
 As neural networks are non-linear with respect
 to their parameters, ensembling all layers---as we do when end-to-end fine-tuning---typically fails, achieving no better accuracy than a randomly initialized neural network \cite{frankle2020linear}. However, as similarly observed by previous work where part of the optimization trajectory is shared \cite{izmailov2018averaging,frankle2020linear, neyshabur2020being}, we find that the zero-shot and fine-tuned models are connected by a linear path in weight-space along which accuracy remains high (explored further in Section~\ref{sec:landscape}).

Remarkably, as we show in Section \ref{sec:results},  WiSE-FT improves accuracy under distribution shift while maintaining high performance on the reference distribution relative to fine-tuned models. 
These improvements come without any additional computational cost as a single set of weights is used.






